%% Section 8: related work. Budget: 1.5 pages (notes/paper-outline.md,
%% section 8, which lists the planned paragraphs and the Pitts draft text).
%% TODO: write the section.

\section{Related Work}
\label{sec:related}

%% Relations for adequacy versus relations between two denotations.
The tower
$\ALR_0, \ALR_1, \ldots$ of the approximant logical relation
(Section~\ref{sec:approx-lr}) has the same shape as the one that
\citet{Reynolds:1974aa} used to relate the direct and continuation
semantics of a language. He built his relation stage by stage along
the inverse-limit construction of the domains
\citep{Scott:1971aa,Smyth:1982aa}. His index ranges over the stages of
both domains, and his relation is the intersection over all stages.
Our index bounds only the depth of the finite target side, and depth
irrelevance (Lemma~\ref{lem:depth-irrelevant}) makes an intersection unnecessary.

%% Graph and filter models; consistency.
Consistency has a precedent in the $\lambda$-calculus model that
\citet{Scott:1982aa} builds from information systems, whose function
entries pair consistent sets of inputs and outputs. As in the coherent
domains of \citet{Plotkin:1978aa}, our consistency is pairwise: a set
is consistent when its elements are pairwise consistent. Our rule for
entries is the consistency of a function space of information systems
\citep{Larsen:1984aa} in its pairwise form
\citep{Schwichtenberg:2006aa}.
